\documentclass[openany,fontset=fandol]{ctexbook}
\usepackage{amsmath,amssymb}
\usepackage[a4paper,margin=25mm]{geometry}
\usepackage[unicode,hidelinks]{hyperref}
\setlength{\emergencystretch}{2em}
\raggedbottom
\newcommand{\PP}{\mathbb P}
\newcommand{\Ysing}{Y^{\mathrm{sing}}}
\newcommand{\Ysup}{Y^{\mathrm{sup\,sing}}}
\newcommand{\Pdual}{\check P}
\DeclareMathOperator{\codim}{codim}
\begin{document}
\title{代数几何基础第五卷预备笔记\\奇异参数的维数与超奇异判据}
\author{第105批\quad 法、英、中文三方校订}
\date{2026年11月24日}
\maketitle

\section*{范围与页缝}
从法稿物理39页尾的承接句续至物理40页末定理8.7(c)起笔；对照英编稿物理27页末、28页以及第16段。查读法41、42和英29以辨接续。法40末的(c)句在法41首不能直接相接；英28也在同处误接。译文保留可辨部分，不把补证当作复原的手稿。法40的定理号可由法41、42回指核作8.7；这一核号不自动解决法36独立段首的覆写号。本批未并入1158页合订本。
 
\section*{奇异截面的参数像}
在这番颇费力的符号推演之前，只看点集也可以先作如下观察。若 $S=\operatorname{Spec} k$，则在非空情形 $\dim\Ysing=r-1$，而 $\dim\Pdual=r$；所以 $\Ysing$ 在 $\Pdual$ 中的像闭包余维至少为一。这又给出(2.12)在这里的结论；所用论证与(2.12)的论证本质上没有区别。\footnote{法40说像，正文用像闭包来定义其余维，不预设像闭。若 $\Ysing$ 为空，则所有参数都避开奇异零；不能给空集赋通常的维数 $r-1$。这里重新得到的是(2.12)的相应方向：光滑处的泛截面光滑且具期望维数；(2.12)的逆向见第89批的平坦下降论证。}

注意，在许多情形下，这个集合的确是余维一（见下文）。因此一般不能指望找到所有成员都几何光滑的超平面截面线性笔。不过往往仍能找到一个没有任何超奇异点的截面笔，因为 $\Ysup$ 在 $\Pdual$ 中的像通常为余维二；避开该像只需要其闭包余维至少为二。\footnote{法40说常见情形为余维二，不是普遍等式；例如坏集为空时也不是余维恰好二。原稿“光滑”在此针对截面具有期望维数的情形。法40边上移入的首行给出后述(b)中源、靶的相对维数。} 下面先归纳这一问题中与微分有关的要点。

\paragraph{译校补证：像不必闭，余维估计仍成立。}
设 $X$ 为有限型。由上批，$\Ysing=\PP_X(\mathcal N_{X/P})$。按 $X$ 的光滑维数分块，在正余法秩的每一块上，$\Ysing$ 的维数都为 $d+(r-d-1)=r-1$。一个不可约分支到其像闭包是支配的，从而诱导函数域嵌入；\href{https://stacks.math.columbia.edu/tag/0A21}{有限型整概形的维数等于函数域超越次数}，故像闭包维数不超过源分支维数。如果 $X/k$ 适当，则 $X\times_k\Pdual\to\Pdual$ 适当，$\Ysing$ 是源中的闭子概形，因而投影本身为适当态射，像确为闭集。这一步不要求 $f$ 是闭浸入。

例如线性嵌入 $X=\PP^d\subsetneq \PP^r$ 且 $d\ge1$ 时，切超平面恰是包含整个 $X$ 的超平面，故参数像为 $\PP^{r-d-1}$，余维是 $d+1$，未必为一。另一方面，若像已经闭且有超曲面分支，则任一几何参数直线都与该分支相交。若闭坏集 $D$ 的余维至少二而 $r\ge2$，过 $D$ 中一点的直线族维数为 $r-1$；故相交直线族维数至多 $(r-2)+(r-1)=2r-3$，小于全直线族的 $2r-2$。可在非空开集中取避开 $D$ 的几何直线；无限底域上也可用 Grassmann 概形的仿射大胞腔取 $k$-直线。有限底域不在此推出一定存在 $k$-直线；$r=1$ 时余维至少二的闭坏集为空。

\newpage
\section*{定理8.7：微分判据}
\paragraph{(a)}
这里由第16段定义的 $\Ysing$，正是前述嵌入 $Y$ 的 $\PP_X(\mathcal N_{X/P})$。这是闭子概形的等同，而非仅为底集等同。

\paragraph{(b)}
$\Ysup$ 的底集，正是以下态射的分歧点集：
\[
 \pi:\Ysing=\PP_X(\mathcal N_{X/P})\longrightarrow\Pdual.
\]
源、靶俱为 $S$-光滑概形，非空处的相对维数分别为 $r-1$ 与 $r$。设 $y$ 在 $\Pdual$ 中的像为 $\xi$；则 $\pi$ 在 $y$ 无分歧，当且仅当 $y$ 对纤维截面 $\varphi_\xi$ 几何上是普通奇异零点。\footnote{这里的分歧是“不无分歧”，不是“不\'etale”；源、靶相对维数本来就不同。参看法40正文及页顶移写行，并与 \href{https://webusers.imj-prg.fr/~leila.schneps/grothendieckcircle/EGA/EGA5trans.pdf}{英编稿的第16段、物理62--65页} 互校。该段判别二次型的极化是否非退化；特征二时不能改用除以2的公式。}

\paragraph{译校补证：纤维的切空间。}
将 $\xi$ 的纤维扩到几何点域 $K$，在 $X_K$ 的相应有理点 $x$ 附近取光滑坐标 $t_1,\ldots,t_d$ 及线丛标架，令 $u$ 是 $\varphi_\xi$ 的局部函数。奇异零纤维由 $u$ 和一阶相对偏导数 $\partial_i u$ 的零式定义。在奇异零点 $x$ 处，$u$ 的线性项为零，因此它不给该点的切空间增加约束。各 $\partial_i u$ 的一次项，正是 $u$ 的二次初始形式 $Q$ 的极化线性映射在各坐标上的分量。具体地，
\[
 B_Q(v,w)=Q(v+w)-Q(v)-Q(w),\qquad
 T_{y}(\pi^{-1}(\xi)_K)=\{v:B_Q(v,w)=0\text{ 对所有 }w\}.
\]
例如在特征二时，平方项的极化为零，交叉项仍计入，整个计算未除以2；单变量 $Q=t^2$ 的极化即为零。由 \href{https://stacks.math.columbia.edu/tag/0B2G}{几何纤维上的切空间判据}，$\pi$ 在该点无分歧恰好等价于这个切空间为零，亦即 $B_Q$ 非退化。无分歧可在几何换底后检测，这就给出了(b)中几何定义与态射判据的对应。

\paragraph{(c) 可辨起笔}
设 $S=\operatorname{Spec} k$，$y\in\Ysing(k)$ 在 $X$ 和 $\Pdual$ 中的投影分别为 $x$ 和 $\xi$；考察由 $d\pi_y$ 的像所确定的、经过 $\xi$ 的射影线性子空间 $H'$。［法40止于切映射所属态射的引出语；(c)的后半原句在现存页缝未找到。］\footnote{英28注40的第一投影可由法40及语法判为 $x$。英28紧接“其像的闭包”及支配态射 $g$，与法41开头相同，不能据此认定它仍属(c)。法41、42已分别回指8.7(b)、8.7(c)、8.8，可知定理条目与后续命题应分开。}

\newpage
\section*{对偶切空间：接续所需的独立补证}
下述推导用于理解法42对8.7(c)的回指，不声称逐字恢复遗失处。令 $a=f(x)$，在 $P\times_k\Pdual$ 中以配对 $\langle a,\xi\rangle=0$ 表示关联条件；写在各自的局部仿射标架中。对来自 $\Ysing$ 的切向量 $(v,w)$ 微分这一关系：
\[
 \langle df_x(v),\xi\rangle+\langle a,w\rangle=0.
\]
第一项为零，因为 $\xi$ 的截面在 $x$ 处的一阶导数为零。因此 $\langle a,w\rangle=0$。也就是说，$\operatorname{im}(d\pi_y)$ 包含于对偶空间中超平面
\[
 a^\perp=\{\xi'\in\Pdual : a\in H_{\xi'}\}
\]
在 $\xi$ 处的切空间。此条件与局部标架的选择无关，因为附随改变的项均乘上在该点已为零的关联函数。

由(b)，若 $y$ 为普通奇异零，则 $d\pi_y$ 为单射；它的像及 $T_\xi(a^\perp)$ 都有维数 $r-1$，故二者相等。切空间与过 $\xi$ 的射影线性子空间的典范对应于是给出 $H'=a^\perp$，进而由 $H'$ 恢复 $a$。反之，若 $H'=a^\perp$，切映射的秩为 $r-1$，从而为单射，回到普通奇异零的情形。这里只能凭此恢复 $a\in P$；若要由 $a$ 唯一恢复 $x\in X$，必须再用 $f$ 为浸入等保证点唯一的条件，不能只用无分歧。

\paragraph{下一步。}
法41以给像闭包赋约化结构起笔，其后可见函数域扩张、双有理、非密切的切超平面以及维数断言；英28把起首若干行误接进8.7(c)。下一批从法41至43校读8.8及8.9，区分几何整性、浸入、底域扩张与适当性分别在哪一步用到；在找不到更好底本时，单列缺文，不把合理推导当成原作者文字。
\end{document}
